\documentclass{article}
\usepackage{amsmath,amssymb,natbib,booktabs,multirow}

% Notation used throughout the paper
\newcommand{\vect}[1]{\mathbf{#1}}
\newcommand{\R}{\mathbb{R}}
\newcommand{\E}{\mathbb{E}}
\DeclareMathOperator*{\argmin}{arg\,min}
\newtheorem{theorem}{Theorem}

\title{Small Models, Careful Data: A Note on Fine-Tuning Budgets}
\author{Ada Example \and Grace Sample}
\date{October 2026}

\begin{document}
\maketitle

\begin{abstract}
We study how far a fixed compute budget goes when it is split between
pre-training and fine-tuning. On three benchmarks, a model with
$10\times$ fewer parameters matches a larger baseline when its fine-tuning
data is filtered first. We give a simple bound that explains when this
happens.
\end{abstract}

\section{Introduction}
\label{sec:intro}

Scaling laws suggest that loss falls predictably with parameters and
data~\citep{kaplan2020scaling}. Yet practitioners rarely train at the
compute-optimal point: they fine-tune what already exists
\citep[see][ch.~3]{howard2018universal}. \citet{devlin2019bert} showed that
a single pre-trained model transfers to many tasks with little task data.

We ask a narrower question: \emph{given a budget, is it better to spend it on
a bigger model or on cleaner fine-tuning data?} Section~\ref{sec:method}
sets up the problem and Table~\ref{tab:results} summarizes the answer.

\section{Method}
\label{sec:method}

Let $\vect{x} \in \R^d$ be an input and $y \in \{1, \dots, K\}$ its label.
We fine-tune parameters $\theta$ by minimizing the empirical risk
\begin{align}
  \hat{\theta} &= \argmin_{\theta} \; \frac{1}{n} \sum_{i=1}^{n}
    \ell\left(f_\theta(\vect{x}_i), y_i\right) + \lambda \|\theta\|_2^2,
    \label{eq:erm} \\
  \ell(z, y) &= -\log \frac{e^{z_y}}{\sum_{k=1}^{K} e^{z_k}}.
\end{align}
Filtering keeps the examples whose loss under a reference model falls
below a threshold $\tau$, which removes about 30\% of the data.

\begin{theorem}[Budget split]
If the label noise rate is $\eta < \frac{1}{2}$, filtering reduces the
excess risk of \eqref{eq:erm} by a factor of at least $1 - 2\eta$.
\end{theorem}

\section{Results}

\begin{table}[t]
  \centering
  \caption{Accuracy (\%) on three benchmarks. Best result in bold.}
  \label{tab:results}
  \begin{tabular}{lccc}
    \toprule
    \multirow{2}{*}{Model} & \multicolumn{3}{c}{Benchmark} \\
    \cmidrule(lr){2-4}
     & A & B & C \\
    \midrule
    Large, raw data      & 81.2 & 74.5 & 68.9 \\
    Small, raw data      & 77.0 & 70.1 & 63.4 \\
    Small, filtered data & \textbf{81.9} & \textbf{75.3} & 68.1 \\
    \bottomrule
  \end{tabular}
\end{table}

The small model with filtered data matches or beats the large one on two of
three benchmarks (Table~\ref{tab:results}), at a tenth of the inference
cost. On benchmark C the gap is within one standard deviation.

\section{Conclusion}

Cleaning the fine-tuning set is a cheaper lever than adding parameters,
as long as the noise rate is high enough for the bound in
Section~\ref{sec:method} to matter.

\bibliographystyle{plainnat}
\bibliography{refs}

\end{document}
